\begin{lemma}[$L^s$ is linear on its domain, and the domain is everything for $s \leq 0$]
Let $s \in \mathbb{R}$, let $\operatorname{Dom}(L^s)$ be as in
\noderef{FracPowDomain} and let $L^s$ be as in \noderef{FracPow}.  Then:
\begin{enumerate}
\item if $s \leq 0$, then every $v \in L^2(\gamma_n)$ lies in
$\operatorname{Dom}(L^s)$;
\item for all $v, w \in \operatorname{Dom}(L^s)$, also
$v + w \in \operatorname{Dom}(L^s)$ and $L^s(v+w) = L^s v + L^s w$;
\item for every $c \in \mathbb{R}$ and $v \in \operatorname{Dom}(L^s)$, also
$c\,v \in \operatorname{Dom}(L^s)$ and $L^s(c\,v) = c\,L^s v$.
\end{enumerate}
\end{lemma}
\begin{proof}
Throughout we use the following three facts, in the form in which they are
stated.

\emph{(a) The coefficient map.}  By \noderef{HermiteBasis} the map
$v \mapsto (v_\alpha)_{\alpha}$, $v_\alpha = \langle v, h_\alpha\rangle$, is an
isometric isomorphism of $L^2(\gamma_n)$ onto $\ell^2(\mathbb{N}_0^n)$.  Being an
isomorphism of real Hilbert spaces it is in particular linear and injective, so
for all $v, w \in L^2(\gamma_n)$ and $c \in \mathbb{R}$
\[
  (v+w)_\alpha = v_\alpha + w_\alpha,
  \qquad (c\,v)_\alpha = c\,v_\alpha
  \qquad (\alpha \in \mathbb{N}_0^n),
\]
every $v \in L^2(\gamma_n)$ has $\sum_\alpha |v_\alpha|^2 = \|v\|_2^2 < \infty$,
and two elements of $L^2(\gamma_n)$ with the same coefficient family are equal.

\emph{(b) The domain criterion.}  By \noderef{FracPowDomain},
$v \in \operatorname{Dom}(L^s)$ holds precisely when the multiplied family
$m_s\bigl((v_\alpha)_\alpha\bigr)$ of \noderef{FracPowCoeff} is square summable,
where by \noderef{FracPowCoeff}
\[
  m_s(c)_\alpha =
  \begin{cases} 0, & \alpha = 0, \\ |\alpha|^s c_\alpha, & \alpha \neq 0 .\end{cases}
\]

\emph{(c) The value.}  By \noderef{FracPow}, for $v \in \operatorname{Dom}(L^s)$
the element $L^s v$ is the unique element of $L^2(\gamma_n)$ whose coefficient
family is $m_s\bigl((v_\alpha)_\alpha\bigr)$.

\medskip
\emph{Proof of (1).}  Let $s \leq 0$ and $v \in L^2(\gamma_n)$.  Fix
$\alpha \neq 0$.  By \noderef{MultiIndexSum} we have $|\alpha| \geq 1$, so
$\log|\alpha| \geq 0$ and hence, since $s \leq 0$,
\[
  |\alpha|^s = e^{s\log|\alpha|} \leq e^0 = 1 ,
  \qquad\text{and also } |\alpha|^s > 0 .
\]
Therefore, by (b),
\[
  \bigl|m_s\bigl((v_\alpha)_\alpha\bigr)_\alpha\bigr|
    = |\alpha|^{s}\,|v_\alpha| \leq |v_\alpha|
  \qquad (\alpha \neq 0),
\]
while for $\alpha = 0$ the left-hand side is $0 \leq |v_0|$.  Thus the family
$m_s((v_\alpha)_\alpha)$ is dominated in absolute value, term by term, by the
family $(v_\alpha)_\alpha$, whence
\[
  \sum_\alpha \bigl|m_s\bigl((v_\alpha)_\alpha\bigr)_\alpha\bigr|^2
    \leq \sum_\alpha |v_\alpha|^2 < \infty
\]
by (a) and the comparison test for series of nonnegative terms.  So
$m_s((v_\alpha)_\alpha)$ is square summable, and $v \in \operatorname{Dom}(L^s)$
by (b).

\emph{Proof of (2).}  Let $v, w \in \operatorname{Dom}(L^s)$.  We first check that
$m_s$ is additive on coefficient families.  Indeed, using
$(v+w)_\alpha = v_\alpha + w_\alpha$ from (a), at $\alpha = 0$ both
$m_s((v+w)_\alpha)_0 = 0$ and
$m_s((v_\alpha))_0 + m_s((w_\alpha))_0 = 0 + 0 = 0$, while for $\alpha \neq 0$
\[
  m_s\bigl(((v+w)_\alpha)_\alpha\bigr)_\alpha
    = |\alpha|^s (v_\alpha + w_\alpha)
    = |\alpha|^s v_\alpha + |\alpha|^s w_\alpha
    = m_s\bigl((v_\alpha)_\alpha\bigr)_\alpha
      + m_s\bigl((w_\alpha)_\alpha\bigr)_\alpha .
\]
Hence
\[
  m_s\bigl(((v+w)_\alpha)_\alpha\bigr)
    = m_s\bigl((v_\alpha)_\alpha\bigr) + m_s\bigl((w_\alpha)_\alpha\bigr)
  \tag{$*$}
\]
as families indexed by $\mathbb{N}_0^n$.  Both summands on the right are square
summable by (b), because $v, w \in \operatorname{Dom}(L^s)$; and a sum of two
square-summable families is square summable, since
$|a+b|^2 \leq 2|a|^2 + 2|b|^2$ termwise.  So the left-hand side of $(*)$ is
square summable, i.e.\ $v + w \in \operatorname{Dom}(L^s)$ by (b).

It remains to identify $L^s(v+w)$.  By (c) the coefficient family of
$L^s(v+w)$ is the left-hand side of $(*)$, and the coefficient families of
$L^s v$ and $L^s w$ are the two summands on the right.  By linearity of the
coefficient map (a), the coefficient family of $L^s v + L^s w$ is the sum of
those two, i.e.\ the right-hand side of $(*)$.  So $L^s(v+w)$ and
$L^s v + L^s w$ are two elements of $L^2(\gamma_n)$ with the same coefficient
family; by the injectivity in (a) they are equal.

\emph{Proof of (3).}  Let $c \in \mathbb{R}$ and $v \in \operatorname{Dom}(L^s)$.
Using $(c\,v)_\alpha = c\,v_\alpha$ from (a), at $\alpha = 0$ both
$m_s(((c v)_\alpha))_0 = 0$ and $c\,m_s((v_\alpha))_0 = c\cdot 0 = 0$, while for
$\alpha \neq 0$
\[
  m_s\bigl(((c\,v)_\alpha)_\alpha\bigr)_\alpha
    = |\alpha|^s (c\,v_\alpha)
    = c\,\bigl(|\alpha|^s v_\alpha\bigr)
    = c\, m_s\bigl((v_\alpha)_\alpha\bigr)_\alpha .
\]
Hence
\[
  m_s\bigl(((c\,v)_\alpha)_\alpha\bigr)
    = c\, m_s\bigl((v_\alpha)_\alpha\bigr)
  \tag{$**$}
\]
as families.  The right-hand side is square summable, because
$\sum_\alpha |c\,x_\alpha|^2 = c^2 \sum_\alpha |x_\alpha|^2$ and
$m_s((v_\alpha)_\alpha)$ is square summable by (b).  So
$c\,v \in \operatorname{Dom}(L^s)$ by (b).  Finally, by (c) the coefficient
family of $L^s(c\,v)$ is the left-hand side of $(**)$, and by (c) together with
the linearity in (a) the coefficient family of $c\,L^s v$ is the right-hand side
of $(**)$; the two agree, so the injectivity in (a) gives
$L^s(c\,v) = c\,L^s v$.
\end{proof}
